\newcommand{\F}{\mathbb{F}}
\newcommand{\C}{\mathbb{C}}
\newcommand{\an}{\mathscr{A}} 
\newcommand{\Q}{\mathbb{Q}}
\renewcommand{\hom}{\operatorname{Hom}}
\newcommand{\ehom}{\operatorname{End}}
\colorlet{gris}{black!60}

\-\vspace{-0.3cm}
  {\hrule \vspace{-0.15cm}
  \centering \textbf{3} \vspace{0.1cm}
  \hrule}

A. As métricas \( d_2, d_1, d_\infty\) definem a mesma topologia em \(\R^n\).   

\textcolor{gris}{\(\bullet \ \ (\tau_2 = \tau_\infty)\). Seja \(|x_j| := \sup_{i\leq n} |x_i|\). 
\[\|\x\|_2^2 = \sum x_i^2 \leq n x_j^2 = n \|\x\|_\infty^2\ \leftrightharpoons \ \|\x\|_2 \leq \sqrt{n}\|\x\|_\infty, \]
logo \(B_\infty(\x; \epsilon)\subset B_2(\x; \sqrt{n}\epsilon)\) e \(\tau_2 \subset \tau_\infty\). Por outro lado, 
\[\|\x\|_\infty^2 = x_j^2 \leq \sum x_i^2 = \|\x\|_2^2 \ \leftrightharpoons \ \|\x\|_\infty \leq \|\x\|_2,  \]
portanto tambem \(B_2(\x; \epsilon ) \subset B_\infty(\x; \epsilon )\) e \(\tau_\infty \subset \tau_2\). \\ 
\(\bullet \ \ (\tau_\infty = \tau_1)\). Observe que, 
\[\|\x\|_\infty = |x_j| \leq \sum |x_i| = \|\x\|_1, \] 
logo \(B_1(\x;\epsilon) \subset B_\infty(\x;\epsilon)\) e \(\tau_\infty \subset \tau_1\). Por outro lado, 
\[\|\x\|_1 = \sum |x_i| \leq n |x_j| = n \|\x\|_\infty, \]
portanto tambem \(B_\infty(\x;\epsilon) \subset B_1(\x;\epsilon)\) e \(\tau_1 \subset \tau_\infty\). 
}

\-\vspace{-0.3cm}
  {\hrule \vspace{-0.15cm}
  \centering \textbf{4} \vspace{0.1cm}
  \hrule}

E. Dado \(A\subset X\) defina a \emph{fronteira} como sendo \(\partial A := \overline{A} \cap \overline{X\setminus A }\). (a) Mostre que \(\overline{A} = A \cup \partial A\). 
\textcolor{gris}{
  \begin{align*}
  A \cup \partial A = A \cup (\overline{A} \cap \overline{X\setminus A }) &= (A \cup \overline{A}) \cap  (A \cup \overline{X \setminus A}) \\ 
  &= \overline{A} \cap X = \overline{A}.
  \end{align*}
}
  (b) Mostre que \(\overset{\  \circ}{A} = A \setminus \partial A\).  
\textcolor{gris}{\begin{align*}
A\setminus \partial A &= A \cap [X \setminus (\overline{A} \cap \overline{X\setminus A})]\\ 
&=A \cap [(X \setminus \overline{A}) \cup (X\setminus \overline{X\setminus A})] \\ 
&=[A \cap (X \setminus \overline{A})] \cup [A \cap \operatorname{int}(X\setminus (X\setminus A))] \\ 
&=[\varnothing] \cup [A \cap \overset{\ \circ }{A}]  = \overset{\ \circ}{A} . 
\end{align*}}

  {\hrule \vspace{-0.15cm}
  \centering \textbf{5} \vspace{0.1cm}
  \hrule}

C. Dados \(f \in C[a,b], \ A \subset [a,b]\) finito e \(r>0\), defina, 
\[V(f, A, r):= \{ g \in C[a,b] : |g(x) - f(x)| < r, \ \ \forall x \in A\}.\]
Mostre que os conjuntos \(V(f, A, r)\) formam uma base de vizinhanças de \(f\) pra uma certa topologia em \(C[a,b]\). 

\textcolor{gris}{Seguindo a \(\square \ 5.7.\) vamos mostrar que a familia verifica as propriedades. (a) Claramente \(\forall U \in \U_f, \ f \in U := V(f,A,r)\) pois \(|f(x) - f(x)| = 0 < r\); (b) Sejam \(U:= V(f, A_1, r_1), \ V := V(f, A_2, r_2) \in \U_f\). Defina \(A:= A_1 \cup A_2, \ r := \min \{r_1, r_2\}\) e \(W:= V(f, A, r) \in \U_f\). Se \(g \in W\), então \(|g(x) - f(x)|< r_1, r_2, \ \forall x \in A_1 \cup A_2\), segue que \(g \in U \cap V\); (c) Sejam \(U = V(f,A,r)\) e \(g \in V:= U\). Defina \(\delta := r - \max_A |g(x)-f(x)| > 0\) e \(W := V(g,A,\delta)\), assim, se \(h \in W\) e \(x\in A\) temos 
\begin{align*}
  |h(x) - f(x) | &\leq |h(x)- g(x)| + |g(x)-f(x)|\\ 
  &< \delta + \max_A |g(x) - f(x)|  = r,  
\end{align*} 
portanto \(h \in U \) e \(W \subset U\).   
}

D. Dado \(A \subset X\), defina o conjunto dos puntos de acumulação \(A' := \{x \in X \ | \  \forall U \in \U_x, \ \exists a \in U\cap A  : a \neq x\}\). Mostre que \(\overline{A} = A \cup A'\).  

\textcolor{gris}{(\(\subseteq \)) Seja \(x \in \overline{A}: x \notin A\), então \(\forall U \in \U_x, \ U\cap A \neq \varnothing \), mas \(x \notin A\), portanto \(\forall U \in \U_x, \ \exists a \neq x : a \in U\cap A \ \leftrightharpoons \ x \in A'\). (\(\supseteq \)) Seja \(x \in A' \setminus A\), então  \(\forall U \in \U_x, \ \exists a \in U \cap A : a \neq x\), em particular \(\forall U \in \U_x, \ U \cap A \neq \varnothing \ \leftrightharpoons \ x \in \overline{A}\).  }

E. De um exemplo de \(A \subset \R : \) todos os seguintes conjuntos sejam diferentes entre si: 
\[A, \ \overline{A}, \ \overset{\ \circ}{A}, \ \overset{ \circ }{\overline{A}},  \ \overline{\overset{ \circ}{A}}, \ \overline{ \overset{ \circ }{\overline{A}} }, \ \overset{\circ }{ \overline{\overset{ \circ}{A}} }.\]
\textcolor{gris}{\[A:= \{3\} \cup (\Q \cap [-1,0]) \cup (0,1) \cup (1,2]\]
\begin{align*}
  \overline{A} = \{3\} \cup [-1,2] \ \ &; \ \ \overset{\ \circ}{A} = (0,1)\cup (1,2) \\ 
\overset{ \circ }{\overline{A}}= (-1,2)\ \  &; \ \ \overline{\overset{ \circ}{A}} = [0,2]\\ 
\overline{ \overset{ \circ }{\overline{A}} } =  [-1,2]\ \ &; \ \ \overset{\circ }{ \overline{\overset{\circ}{A}}  }= (0,2)
\end{align*}
}

  {\hrule \vspace{-0.15cm}
  \centering \textbf{6} \vspace{0.1cm}
  \hrule}

A. Prove que o segundo axioma de enumerabilidad implica o primeiro.  

\textcolor{gris}{Sendo \(\mathcal{B}\) uma base enumeravél pra \(X\) defina \(\mathcal{B}_x := \{V \in \mathcal{B} : x \in V\}\). Dada \(U \in \U_x\), temos que \(x \in \operatorname{int}U=: U_0\ab\subset X\), logo \(\exists \mathcal{C} \subset \mathcal{B} : U_0 = \bigcup \{V : V \in \mathcal{C}\}\). Em particular, \(\exists V \in \mathcal{B} : x \in V \subset U_0 \subset U\). Assim,  \(V \in \mathcal{B}_x\) é tal que \(V \subset U\), concluindo que \(\mathcal{B}_x\) é base de vizinhanças, enumeravél pois \(|\mathcal{B}_x| \leq |\mathcal{B}|\).   } 

C. Prove que os intervalos \([a,b) \) com \(a<b\) formam uma base pra uma topologia \(\tau_\ell \) em \(\R\) mais fina que \(\tau_2\). 

\textcolor{gris}{Vamos verificar a propriedades na \(\square \ 6.6.\). (a) Dado \(x \in \R\) temos \(x \in [x, x+1)\), logo \(x \in \bigcup \{[a,b) : a <b\}  = \R\); (b) Dado \(x \in [a,b) \cap [c,d)\) temos, necessariamente, 
\[  \max\{a,c \} =: e \leq x < f := \min \{b,d\}, \]
logo, \(x\in W := [e,f) \subset [a,b) \cap [c,d)\), onde \(W \in \mathcal{B}\). Segue que \(\mathcal{B}\) é base de \((\R, \tau_\ell) =: \R_\ell\) (\emph{reta de Sorgenfrey}). \\
}

  {\hrule \vspace{-0.15cm}
  \centering \textbf{8} \vspace{0.1cm}
  \hrule}

B. Prove que uma função \(f : X \to Y \) é contínua \(\leftrightharpoons\) \(f(\overline{A}) \subset \overline{f(A)}\), pra cada \(A \subset X\).  

\textcolor{gris}{(\(\Rightarrow\)) Sejam \( f(x) = y \in f(\overline{A}): x \in \overline{A}\) e \(V\in \U_y\), pela continuidade \(x \in f^{-1}(V) \in \U_x\). Como \(x \in \overline{A}\), \(f^{-1}(V) \cap A \neq \varnothing\), segue que, 
\[f(f^{-1}(V) \cap A) = f(f^{-1}(V)) \cap f(A) \subset V \cap f(A) \neq \varnothing, \]
isso \(\forall V \in \U_y\), logo \(y \in \overline{f(A)}\). \\ 
(\(\Leftarrow\))  Sejam \(F\ce\subset Y\) e \(A := f^{-1}(F)\). Note que, 
\[f(A)\subset f(\overline{A}) \subset \overline{f(A)} = \overline{f(f^{-1}(F))} \subset \overline{F}= F, \]
assim, \(\overline{A} \subset f^{-1}(f(\overline{A})) \subset f^{-1}(F) = A\ce \subset X\). 
}

D. Prove que se \(f,g : X \to \R\) são contínuas em \(a \in X\), então as funções \(f+g\) e \(fg\) são tambem contínuas em \(a\).  

\textcolor{gris}{(\(+\)) Sejam \(\epsilon > 0\) e \( U_1, U_2  \in \U_a\) tais que: 
\[f(U_1)\subset B(f(a), \nicefrac{\epsilon}{2}) \text{ \ \ e \ \ }g(U_2)\subset B(g(a), \nicefrac{\epsilon}{2}).\]
Dado \(x \in U:= U_1 \cap U_2 \in \U_a\) temos, 
\begin{align*}
  |(f+g)(x) - (f+g)(a)| &\leq  |f(x) -f(a)| + |g(x)-g(a)| \\ 
& < \nicefrac{\epsilon }{2} + \nicefrac{\epsilon }{2} = \epsilon, 
\end{align*}
logo \((f+g)(U) \subset B((f+g)(a), \epsilon)\), \(f+g\) é contínua em \(a\).   \\ 
(\(\cdot\)) Primeiro note que, 
\begin{align*}
  |(fg)(x) - (fg)(a)| &= |(fg)(x) + f(x)g(a) \\ 
  & \hspace{1.9cm} - f(x)g(a) - (fg)(a)|\\ 
  &\leq |f(x)||g(x)-g(a)| \\ 
  & \hspace{1.9cm}+ |g(a)||f(x)-f(a)| 
\end{align*} 
\(\bullet \ \exists U_1 \in \U_a : |f(x)-f(a)| < 1 \), em particular temos \(|f(x)| < 1 + |f(a)|\).  \\ 
\(\bullet \ \exists U_2 \in \U_a : |g(x)-g(a)| < \frac{\epsilon}{2(1 + |f(a)|)} \). \\ 
\(\bullet \ \exists U_3 \in \U_a : |f(x)-f(a)| < \frac{\epsilon}{2(1 + |g(a)|)} \). \vspace{0.3cm}\\ 
Finalmente, sendo \(x \in U := U_1 \cap U_2 \cap U_3 \in \U_x\) temos na equação acima, 
\[ |(fg)(x) - (fg)(a)| < \nicefrac{\epsilon}{2} + \nicefrac{\epsilon}{2} = \epsilon.\]
Assim, \(fg(U) \subset B((fg)(a);\epsilon)\) e \(fg\) é contínua em \(a\). 
}

\vspace{0.3cm}
  {\hrule \vspace{-0.15cm}
  \centering \textbf{9} \vspace{0.1cm}
  \hrule}

\vspace{0.3cm}
  {\hrule \vspace{-0.15cm}
  \centering \textbf{12+} \vspace{0.1cm}
  \hrule}

\vspace{0.4cm}
%\vspace{0.3cm}
%  {\hrule \vspace{-0.15cm}
%  \centering \textbf{13} \vspace{0.1cm}
%  \hrule}

%\vspace{0.3cm}
%  {\hrule \vspace{-0.15cm}
%  \centering \textbf{15} \vspace{0.1cm}
%  \hrule}


%A. Se \((x_\lambda)_{\lambda \in \Lambda } \subset X\) é tal que \(x_\lambda \to x \), então qualquer subrede \((x_{\phi(\mu)})\) tambem converge a \(x\).  \\ 
%\textcolor{gris}{Sejam \(x\circ \phi : M \to X\) subrede e \(U\in \U_x\) quaisquer, como \(x_\lambda \to x, \ \exists \lambda_0 : (x_\lambda)\subset U\) se \(\lambda \geq \lambda_0\). Como \(\phi \) é cofinal \(\exists \mu_0 \in M : \phi(\mu_0) \geq \lambda_0\), além disso é crescente, portanto \((x_{\phi(\mu)}) \subset U\) se \(\mu \geq \mu_0\) pois \(\phi(\mu) \geq \phi(\mu_0) \geq \lambda_0\). Para \(\phi\) é \(U\) arbitrarias, concluimos então que \(x_{\phi(\mu)} \to x\). 
%}

%B. Se \(x \) é ponto de acumulação de uma subrede de \((x_\lambda)_{\lambda \in \Lambda}\), então \(x\) tambem é ponto de acumulação da propia \((x_\lambda)\). \\ 
%\textcolor{gris}{Dados \(U\in \U_x\) e \(\lambda_0 \in \Lambda\), \(\exists \mu_0 : \phi(\mu_0) \geq \lambda_0\), seguido  \(\exists \mu \geq \mu_0: x_{\phi(\mu)} \in U\). Em particular, \(\lambda :=\phi(\mu)\geq \phi(\mu_0) \geq \lambda_0 \) é tal que \( x_\lambda \in U\), segue que \(x\) é ponto de acumulação da \((x_\lambda)\). }
